\section*{Rahmen und Methode}
\addcontentsline{toc}{section}{Rahmen und Methode}
\markright{Rahmen und Methode}
\label{sec:methode}

Daimler Truck Holding AG, Leinfelden-Echterdingen, ist ein Nutzfahrzeughersteller
(Lkw und Busse), im Dezember 2021 von der damaligen Daimler AG (heute Mercedes-Benz
Group AG) abgespalten. Prim\"arquellen sind die vier bislang vorliegenden
Gesch\"aftsberichte 2022 bis 2025, jeweils englisch, heruntergeladen von
daimlertruck.com. Jede Seitenangabe bezeichnet die \emph{gedruckte} Seite. Der Bericht
entscheidet \"uber den Kauf der Aktie zum Kurs von \je{\Kurs}{EUR} (\KursDatum,
Sekund\"arquelle Yahoo Finance).

\subsection*{Der Unterschied zu GEA: zwei Gesch\"afte in einem Unternehmen}

Anders als bei einem Maschinenbauer besteht Daimler Truck aus einem Industriegesch\"aft
(Lkw- und Busbau) und einem Finanzdienstleistungssegment (Absatzfinanzierung und
Leasing f\"ur die eigenen Kunden). Das Finanzdienstleistungsgesch\"aft f\"uhrt zu
wachsenden Forderungen, deren Ver\"anderung im operativen Mittelzufluss der
\emph{gesamten} Gruppe steht und diesen von Jahr zu Jahr um mehrere Milliarden Euro
verschiebt, ohne etwas \"uber das Industriegesch\"aft auszusagen.

\annahme{Der Bericht rechnet deshalb nicht mit dem operativen Mittelzufluss der Gruppe,
sondern mit dem von Daimler Truck selbst ausgewiesenen \enquote{Free cash flow of the
Industrial Business}: dem freien Zufluss des Lkw- und Busgesch\"afts allein, ohne
Financial Services. Das ist eine \"Ubernahme der eigenen Segmentrechnung des
Unternehmens, keine eigene Herleitung; der Wert erscheint auf derselben Kennzahlenseite
wie Umsatz, EBIT und ROCE und ist gegen die Kapitalflussrechnung nicht rekonstruierbar,
weil das Unternehmen die Aufteilung nicht Zeile f\"ur Zeile offenlegt.}

\annahme{Alle Rechnungen sind Vorsteuerrechnungen auf Ebene des Unternehmens, mit
konstant fortgeschriebenem Zufluss und dem Buchwert des Eigenkapitals als Endwert,
identisch zur Methode der \"ubrigen Berichte dieses Verzeichnisses.}

\subsection*{H\"urde}

\Pz{\Huerde} p.\,a., die Rendite des MSCI World (Gross Returns, USD) \"uber zehn Jahre
nach dem Factsheet zum 31.08.2026\QW{MSCI World Index, Index Factsheet
31.08.2026}{quelle:msci}{19.09.2026}. Als Sensitivit\"at dient \Pz{\HuerdeLang}, die
Rendite seit Auflage des Index (1987).

\subsection*{Nicht verf\"ugbare Angaben}

Die Organgesch\"afte, die \"uber die offizielle Meldeseite des Unternehmens auffindbar
waren, betreffen ausschlie{\ss}lich den Erwerb von Aktien im Rahmen eines
\enquote{Phantom Share Plan} (Vergütungsprogramm), nicht freie Marktk\"aufe; Preis und
Volumen wurden auf der abgerufenen Seite nicht angegeben. Ein solcher planm\"a{\ss}iger
Erwerb ist kein Signal \"uber die Einsch\"atzung des Kurses und wird in diesem Bericht
deshalb NICHT wie die freien Organk\"aufe bei GEA gewertet; das vollst\"andige Register
(Unternehmensregister bzw.\ BaFin-Datenbank) wurde nicht abgerufen. H\"ochst- und
Tiefstkurse der Analystenziele liegen nicht vor, nur der Mittelwert.
